\scanpage{36}
est cependant possible que l'injectivité de (*) reste valable sans hyp. de caractéristique -- à examiner de plus près...

Revenant aux conditions du cor. 1, b), en car. $p>0$, en travaillant avec les groupes fondamentaux premiers à $p$, on trouve que le $a$ dans 1°) est un point \inserted{de $G(K)$} d'ordre une puissance de $p$. Si $G$ est un tore (il suffit pour cela que les ``fibres'' dans la variété élémentaire $Y$ soient de genre 0) on en conclut encore que $a=0$, i.e. $f=g$. Dans le cas général, la démonstration arithmétique prouve \inserted{(\unsure{on verra|on fera})} ceci : si $x\in X$, \deleted{et} \inserted{d'où} $f(x)\inserted{=P}$, $g(x)=Q$ dans $Y$, non seulement $Q-P$ est un élément de $p$-torsion de la jacobienne gén. $G_Y$ de $Y$, pour tt rev. \inserted{étale} galoisien fini $Y'$ de $X$, \inserted{de degré premier à $p$,} pointé par un $P'$ au dessus de $P$, il existe $Q'$ au dessus de $Q$ tel que $Q'-P'$ soit un élément de \unsure{?} $p$-torsion de $G_{Y'}$.

\deleted{Passant aux revêtements par \inserted{rev. abéliens, \unsure{multiples},} isogénies ${}_nG$ ($n$ premier à \unsure{$G$|$p$}), il existe un $Q_n\in G$ tel que $\unsure{{}_n nQ_n|nQ_n}=Q$, et que $Q_n$}
Il se pourrait qu'on puisse en déduire $P=Q$.

\margin{ceci montre que le th. 2 reste vrai, sur un corps de base de car. $p>0$, en travaillant avec les groupes fondamentaux premiers à $p$ -- car là dans la réduction à un énoncé géométrique, \unsure{on voit qu'à des étales près}, on peut supposer que $Y$ est un \uline{produit} de courbes de genre zéro. Dans ce cas d'ailleurs, on est ramené au cas où $Y$ est $=$ une telle courbe,....}
\margin{35}
